\questao{96}{corC}{Ondas e óptica}{Velocidade das ondas: o eco do radar}{Com conta}
\begin{enunciado}
% gerado por ferramentas/md2tex_enunciado.py a partir de questoes/q096.md (não editar à mão)
Uma instalação de radar, usada para rastrear aviões, emite ondas eletromagnéticas com velocidade estimada de 300 mil quilômetros por segundo. Um avião que saiu da rota foi localizado em um raio de 75 quilômetros do aeroporto onde fica o radar. O operador usou o tempo dispendido para enviar e receber a resposta com a finalidade de estimar a localização do avião.\par
\figq[0.28\linewidth]{q096_fig1.png}
\fonte{Disponível em: www.pasternack.com. Acesso em: 1 abr. 2026 (adaptado).}
O tempo, em segundo, medido foi de\par
\begin{altlinha}
\alth{A}{$2{,}0 \times 10^{3}$\,s.}%
\alth{B}{$4{,}0 \times 10^{0}$\,s.}%
\alth{C}{$5{,}0 \times 10^{-1}$\,s.}%
\alth{D}{$1{,}5 \times 10^{-3}$\,s.}%
\alth{E}{$5{,}0 \times 10^{-4}$\,s.}
\end{altlinha}

\end{enunciado}

\quer{quanto tempo a onda do radar leva para ir até o avião e voltar.}

\begin{ideia}
O radar funciona como um \textbf{eco}. Quando você grita na frente de um paredão, o som vai até a parede, bate e volta.
Você só escuta o eco depois que o som fez \textbf{a ida e a volta}.

Com o radar é igual, só que a onda é eletromagnética (da mesma ``família'' da luz) e muito mais rápida: 300 mil quilômetros
em um único segundo. A antena solta a onda, ela bate no avião e retorna. O tempo medido é o da viagem completa.

Para calcular tempo de viagem, a relação é a mesma de um carro na estrada:
\[ \text{tempo} = \frac{\text{distância}}{\text{velocidade}} \]

\begin{center}
\begin{tikzpicture}[font=\small]
  \node[draw=cinza,fill=cinzaclaro,rounded corners=1mm,minimum width=14mm,minimum height=8mm] (r) at (0,0) {radar};
  \node[draw=cinza,fill=cinzaclaro,rounded corners=1mm,minimum width=14mm,minimum height=8mm] (a) at (8,0) {avião};
  \draw[seta,corC] ([yshift=2mm]r.east) -- ([yshift=2mm]a.west) node[midway,above]{ida: 75\un{km}};
  \draw[seta,corB] ([yshift=-2mm]a.west) -- ([yshift=-2mm]r.east) node[midway,below]{volta: 75\un{km}};
\end{tikzpicture}
\end{center}
\end{ideia}

\passoapasso
\npasso{Calcule a distância total percorrida pela onda.}
O enunciado diz que o avião foi localizado ``em um raio de 75 quilômetros''. Entenda assim: o avião está a 75\un{km} do radar.
A onda vai e volta:
\begin{conta}
$d = 75 + 75 = 150\un{km}$
\end{conta}

\npasso{Escreva a velocidade com todos os zeros e confira as unidades.}
O enunciado escreve a velocidade por extenso: ``300 mil quilômetros por segundo''. Em número, é $300\,000\un{km/s}$ (300 seguido de três zeros).
A distância está em \textbf{km} e a velocidade em \textbf{km por segundo}. As duas usam quilômetro, então não é preciso
converter nada. O tempo vai sair em segundos, como a questão pede.

\npasso{Divida a distância pela velocidade.}
\begin{conta}
$t = \dfrac{150}{300\,000}$
\end{conta}
Para simplificar, corte um zero em cima e um embaixo. Depois divida os dois por 15 ($15 \div 15 = 1$ e $30\,000 \div 15 = 2\,000$, porque $30 \div 15 = 2$):
\begin{conta}
$\dfrac{150}{300\,000} = \dfrac{15}{30\,000} = \dfrac{1}{2\,000} = 0{,}0005\un{s}$
\end{conta}
(Para fazer $1 \div 2\,000$: dividir por 2\,000 é dividir por 2 e depois por 1\,000. Metade de 1 é 0,5. Dividir 0,5 por 1\,000 é andar a vírgula três casas para a esquerda: 0,0005.)

\npasso{Escreva em notação científica.}
As alternativas estão com potências de 10. No número 0,0005, conte quantas casas a vírgula anda para a direita até parar
depois do 5: \textbf{4 casas}. Esse é o expoente. Como o número é menor que 1, o expoente é negativo:
\begin{conta}
$0{,}0005 = 5{,}0 \times 10^{-4}\un{s}$
\end{conta}
É meio milésimo de segundo. Faz sentido ser um tempo minúsculo: a onda é rapidíssima.

\resposta{E}{A onda percorre 150\un{km} (ida e volta) a 300\,000\un{km/s}. Isso leva $5{,}0 \times 10^{-4}\un{s}$.}

\begin{outras}
\alt{A} $2{,}0 \times 10^{3}$ aparece quando se faz a divisão ao contrário ($300\,000 \div 150 = 2\,000$). Seriam 2\,000 segundos, mais de meia hora!\par
\alt{C} $5{,}0 \times 10^{-1}$ aparece para quem esquece o ``mil'' e faz $150 \div 300 = 0{,}5$. Seria meio segundo: tempo demais para uma onda tão rápida.\par
\alt{B} $4{,}0 \times 10^{0}$ aparece para quem faz $300 \div 75 = 4$: esqueceu o ``mil'', esqueceu a volta e inverteu a divisão.\par
\alt{D} $1{,}5 \times 10^{-3}$ não sai de nenhuma conta correta com esses dados.
\end{outras}

\begin{pegadinha}[Onde muita gente erra]
\begin{itemize}
  \item \textbf{Esquecer a volta.} Quem usa só 75\un{km} encontra $2{,}5 \times 10^{-4}\un{s}$, que nem aparece nas alternativas.
        Se a sua conta não bate com nenhuma alternativa, confira se usou ida \textbf{e} volta.
  \item \textbf{Contar os zeros em vez das casas.} 0,0005 tem três zeros depois da vírgula, mas a vírgula anda \emph{quatro} casas: $10^{-4}$.
\end{itemize}
\end{pegadinha}
